\begin{lemma}[Order-nine structural configuration; core of Scheide Proposition 4.5]
Let $G$ be a finite loopless multigraph, let $k$ be a natural number, and
suppose that $G$ is critical,
\[
 \chi'(G)=k+1,\qquad \Delta(G)+1\leq k,
\]
and every vertex of $G$ is incident with an edge.  If $G$ is not
elementary and $t(G)<11$, then there are an edge identity $e_0$, a proper
partial $k$-edge-coloring $\varphi$ leaving exactly $e_0$ uncolored, and a
tree-sequence $T$ such that $(T,e_0,\varphi)$ carries a Scheide structural
configuration.
\end{lemma}
\begin{proof}
By \noderef{GlobalMaximumTripleExists}, choose a globally maximum triple
$(T_0,e_0,\varphi_0)$.  Apply \noderef{ScheideBalancedExists} and replace
it by balanced normal data $(T,e,\varphi,N)$.  Write $t=t(G)$ and
$h=h(G)$.  The global field is unchanged, and the balanced condition in
\noderef{BalancedGlobalMaximumTriple} gives $N.\mathsf{height}=h$.

We first determine both numerical parameters.  Since $G$ is assumed
non-elementary, the contrapositive of
\noderef{ScheideSmallHeightElementary} gives $5\leq h$, and the
contrapositive of \noderef{ScheideLargeHeightElementary} gives
$h\leq t-4$.  The initial globally maximum triple and
\noderef{ScheideGraphHeightOdd} show that $h$ is odd.  The maximum-order
assertion \noderef{ScheideGlobalMaximumOrder}, applied to the balanced
triple, shows that $t\geq3$ and is odd.  Together with $t<11$, the inequalities
$5\leq h\leq t-4$ exclude the odd values $t=3,5,7$ and leave only
\[
 t=9,\qquad h=5.                                      \tag{1}
\]
The equality in \noderef{GlobalMaximumTriple} is an equality
$|V(T)|=t$, where $V(T)$ means \noderef{TreeSequenceVertexSet}.  Vertex
injectivity in \noderef{TreeSequence} separately gives
$|V(T)|=T.\mathsf{length}+1$.  Thus (1) says
$T.\mathsf{length}=8$; it is not an assertion that the alternating
vertex--edge sequence has length nine.

Name its individual entries, using the paper's one-based edge labels, as
\[
 T=(y_0,e_1,y_1,e_2,y_2,e_3,y_3,e_4,y_4,
       e_5,y_5,e_6,y_6,e_7,y_7,e_8,y_8).             \tag{2}
\]
Thus $e_i$ is Lean edge index $i-1$; in particular the two balanced
tail pairs $e_5,e_6$ and $e_7,e_8$ are Lean indices $4,5$ and $6,7$.
By \noderef{BalancedGlobalMaximumTriple} there are colors
$\gamma_1,\gamma_2$ with
\[
 \varphi(e_5)=\varphi(e_6)=\gamma_1,
 \qquad
 \varphi(e_7)=\varphi(e_8)=\gamma_2.                \tag{3}
\]
The values are actual colors because \noderef{GlobalMaximumTriple}
contains \noderef{KTriple}, hence \noderef{IsOneEdgeUncolored}; among the
tree identities only $e_1=e$ is uncolored.  Applying
\noderef{TashkinovTree} to $e_5$ and $e_7$ gives, respectively, a vertex
among $y_0,\ldots,y_4$ missing $\gamma_1$ and a vertex among
$y_0,\ldots,y_6$ missing $\gamma_2$.

Let $y_j$ be the first witness missing $\gamma_1$.  If $j=0$, retain the
triple; this is the no-rotation case.  If $j>0$, apply
\noderef{ScheideTrunkRotationPreserves} at the positive trunk index $j$.
The exact relation \noderef{TrunkRotation} cyclically moves $y_j$ to the
first position, changes colors only on the closing trunk cycle, and keeps
every post-trunk vertex and edge identity in its old position.  The
identities $e_5,\ldots,e_8$ have outside-trunk endpoints and so are not
closing-cycle identities.  Their values in (3) are therefore unchanged.
It remains to check missingness at the new first vertex rather than assume
it.  The elementary clause of \noderef{ScheideGlobalMaximumOrder}, read
through \noderef{ColoredVertexSetElementary}, makes a missing color have
at most one witness in $V(T)$.  If $\gamma_1=\alpha$, (4)'s original
$\alpha$-witness forces $j=0$, so this case was not rotated.  If
$\gamma_1=\beta$, its witness forces $j=1$; the shift prescribed by
\noderef{TrunkRotation} exchanges the uncolored identity and the adjacent
$\alpha$-identity, neither of which has color $\beta$, so $\beta$ remains
missing at $y_1$.  If $\gamma_1$ is neither trunk color, every changed
cycle identity has one of the two trunk colors and hence no incident
$\gamma_1$-value changes.  Thus $\gamma_1$ is missing at the new first
vertex in every case.  Relabel
the rotated trunk by $y_0,\ldots,y_4$ and call its two normal colors
$\alpha,\beta$.  We have proved, in both cases, that the current triple is
balanced and globally maximum of height five, that (2)--(3) still hold,
and that
\[
 \gamma_1\in\overline\varphi(y_0),\qquad
 \alpha\in\overline\varphi(y_0),\qquad
 \beta\in\overline\varphi(y_1).                    \tag{4}
\]

The edge $e_5$ joins its new vertex $y_5$ to the five-vertex trunk by the
attachment field of \noderef{TreeSequence}.  Hence it supplies the
specified crossing identity needed by
\noderef{ScheideHeightFiveGammaEdges}.  That lemma yields identities
$f_1,f_2,f_3$ and vertices $v_1,v_2$ outside the trunk, all in $V(T)$,
such that
\[
 f_1=y_1v_1,\qquad f_2=y_4v_2,\qquad f_3=y_2y_3,
 \qquad
 \varphi(f_1)=\varphi(f_2)=\varphi(f_3)=\gamma_1.   \tag{5}
\]
Here and below an equation such as $f_1=y_1v_1$ abbreviates the exact
unordered-endpoint equality; $f_1$ remains an edge identity.

We identify the first balanced pair without collapsing parallel edges.
The attachment of $e_5$ lies on the trunk.  The attachment of $e_6$ also
lies on the trunk: it cannot be $y_5$, since the distinct identities
$e_5,e_6$ would then be adjacent in the sense of
\noderef{AdjacentEdges} and have the same color, contrary to
\noderef{IsProperEdgeColoring}.  Neither identity can attach at $y_0$,
where $\gamma_1$ is missing by \noderef{MissingColors}; equivalently an
incident identity of that color would put $\gamma_1$ in
\noderef{PresentColors}.  Neither can attach at $y_2$ or $y_3$, since the
identity $f_3$ in (5) already supplies the unique incident
$\gamma_1$-edge there.  Thus their attachment vertices are $y_1,y_4$,
one each: sharing one would again make the two distinct identities
adjacent and equally colored.  Properness makes the incident
$\gamma_1$-identity at $y_1$ unique, so the edge attached there is
exactly $f_1$, not merely parallel to it; its other endpoint is $v_1$.
The same argument at $y_4$ identifies the other identity with $f_2$ and
its other endpoint with $v_2$.  Interchange the two complete
edge--new-vertex entries if necessary.  All vertices and edges remain
injective, both attachments remain to earlier trunk vertices, all
\noderef{TashkinovTree} witnesses and the normal trunk remain unchanged,
so the reordered tree is still balanced and globally maximum.  We may
therefore assume
\[
 e_5=f_1,\ y_5=v_1,\qquad e_6=f_2,\ y_6=v_2.       \tag{6}
\]
Apply \noderef{ScheideHeightFiveOutsideLinks} to (5).  It supplies two
different edge identities $a_{56},b_{56}$ with the same endpoint pair
$\{y_5,y_6\}$ and respective colors $\alpha,\beta$.

We next prove the required color distinctness.  Along the normal
height-five closing path, properness and (4) give $\alpha$-edges
$y_1y_2,y_3y_4$ and $\beta$-edges $y_2y_3,y_4y_0$.  Together with
$a_{56}$ and $b_{56}$, these are matchings of $y_0,\ldots,y_6$ which
leave only $y_0$ unmatched in color $\alpha$ and only $y_1$ unmatched in
color $\beta$.  Equations (5)--(6) similarly give a
$\gamma_1$-matching $y_1y_5,y_2y_3,y_4y_6$, leaving only $y_0$
unmatched.  These assertions use \noderef{Incident} to turn endpoint
equalities into incidences and \noderef{IsProperEdgeColoring} to make
each colored incidence unique.  If $\gamma_2$ equaled any of
$\alpha,\beta,\gamma_1$, each attachment of the same-colored identities
$e_7,e_8$ to the earlier seven vertices would either meet a vertex
already incident with that color, contradicting properness through
\noderef{AdjacentEdges}, or meet the unique missing vertex, contradicting
\noderef{MissingColors} and \noderef{PresentColors}.  Hence
\[
 \gamma_2\notin\{\alpha,\beta,\gamma_1\}.           \tag{7}
\]

We also record the other separations that will be used below.  The identity
$e_5=f_1$ is incident with $y_1$ and has color $\gamma_1$, whereas
$\beta$ is missing at $y_1$ by (4).  If $\beta=\gamma_1$, this incidence
would put $\beta$ in \noderef{PresentColors} at $y_1$, contradicting the
definition of \noderef{MissingColors}.  Thus $\beta\ne\gamma_1$.  The
normal path gives $\varphi(e_2)=\alpha$ and
$\operatorname{ends}(e_2)=\{y_1,y_2\}$.  The identities $e_2$ and $e_5$
are distinct by edge injectivity in \noderef{TreeSequence}, and both are
incident with $y_1$ by \noderef{Incident}.  Hence they are adjacent in the
sense of \noderef{AdjacentEdges}; if $\alpha=\gamma_1$, they would be two
distinct adjacent edges with the same color, contrary to
\noderef{IsProperEdgeColoring}.  Therefore
\[
 \beta\ne\gamma_1,\qquad \alpha\ne\gamma_1.          \tag{7a}
\]
Finally $\alpha\ne\beta$: the distinct vertices $y_0,y_1$ miss
$\alpha,\beta$, respectively, and elementarity from
\noderef{ScheideGlobalMaximumOrder}, expanded through
\noderef{ColoredVertexSetElementary}, makes their missing-color sets
disjoint.  Thus $\alpha,\beta,\gamma_1$ are pairwise distinct.

The last two vertices are joined by three separately named identities.
Indeed, \noderef{ScheideGlobalMaximumOrder} says that $V(T)$ is elementary
and closed in the senses of \noderef{ColoredVertexSetElementary} and
\noderef{ColoredVertexSetClosed}.  For each
$c\in\{\gamma_1,\alpha,\beta\}$ its displayed
matching on the first seven vertices has exactly one missing vertex.
Elementarity therefore makes $c$ present at both $y_7,y_8$, and
\noderef{PresentColors} supplies an incident $c$-colored identity $g_c$
at $y_7$.  We now identify its other endpoint rather than suppressing that
step.  By \noderef{Incident}, write
$\operatorname{ends}(g_c)=\{y_7,w_c\}$; the loopless law in
\noderef{LooplessMultigraph} gives $w_c\ne y_7$.  If
$w_c\notin V(T)$, this exact endpoint equality, together with
$y_7\in V(T)$, puts $g_c$ in \noderef{BoundaryEdges} of $V(T)$.  But $c$
is missing at its displayed unmatched vertex in $V(T)$, so closedness in
\noderef{ColoredVertexSetClosed} forbids a $c$-colored boundary identity.
Consequently $w_c\in V(T)$.

The equality defining \noderef{TreeSequenceVertexSet} and (2) say that
$V(T)=\{y_0,\ldots,y_8\}$.  For each $0\le i\le6$, the vertex $y_i$ is
either the unique displayed vertex missing $c$, in which case no
$c$-colored identity is incident there by \noderef{MissingColors}, or is
incident with its displayed matching identity of color $c$.  In the second
case, if $w_c=y_i$, properness would force $g_c$ to equal that unique
displayed identity; its exact endpoint pair lies wholly among
$y_0,\ldots,y_6$, however, whereas $g_c$ is incident with the distinct
vertex $y_7$.  This is impossible.  Thus $w_c$ is none of
$y_0,\ldots,y_6$, and it is not $y_7$ by looplessness.  The displayed
enumeration of $V(T)$ forces $w_c=y_8$.  Repeating this argument for
$c=\gamma_1,\alpha,\beta$ gives identities with exact endpoint pair
$\{y_7,y_8\}$; denote them by
$g_{\gamma_1},g_\alpha,g_\beta$.  They are pairwise distinct because
their colors are pairwise distinct by (7a) and the preceding
$\alpha\ne\beta$ argument.  Parallel identity is retained even though all
three have endpoint pair $\{y_7,y_8\}$.

The $\gamma_2$-witness preceding $e_7$ lies either on the trunk or in
$\{y_5,y_6\}$.  In the first case keep the present data.  In the second,
let the witness be one of $v_1,v_2$ and replace (2) by
\[
 (y_0,e_1,y_1,f_1,v_1,b_{56},v_2,f_2,y_4,
      e_2,y_2,e_4,y_3,e_7,y_7,e_8,y_8).             \tag{8}
\]
We verify this rebasing before relabeling.  The nine vertices in (8) are
exactly the old nine vertices, so vertex injectivity is preserved.  Its
eight edges are distinct identities: $e_1,f_1,f_2,e_2,e_4,e_7,e_8$ are
distinct old tree identities, while $b_{56}$ is the separate
$\beta$-identity between $v_1,v_2$ and differs from the two
$\gamma_1$-identities and from every edge with a different endpoint
pair.  Their endpoint equalities show successively that every new vertex
attaches to an earlier one, and the first identity remains the uniquely
uncolored $e_1=e$; hence (8) satisfies every field of
\noderef{TreeSequence}.

Every Tashkinov witness is also explicit.  The colors of $f_1,f_2$ are
$\gamma_1$, missing at $y_0$; $b_{56}$ has color $\beta$, missing at
$y_1$; $e_2,e_4$ have color $\alpha$, missing at $y_0$; and the
$\gamma_2$ witness in $\{v_1,v_2\}$ precedes both $e_7,e_8$ in (8).
Thus (8) is a \noderef{TashkinovTree}.  The coloring and distinguished
identity did not change, so its \noderef{KTriple} is the same, and its
vertex set still has cardinality $t=9$; hence
\noderef{GlobalMaximumTriple} holds.

Moreover
\[
 y_1,f_1,v_1,b_{56},v_2,f_2,y_4,q,y_0
\]
is a simple alternating $(\gamma_1,\beta)$ incidence path with exactly
the first five vertices of (8), where $q$ is the old closing identity.
It therefore supplies \noderef{NormalGlobalMaximumTriple} of height five.
The remaining pairs $e_2,e_4$ and $e_7,e_8$ have equal colors
$\alpha$ and $\gamma_2$, respectively, so
\noderef{BalancedGlobalMaximumTriple} holds.  Only now relabel positions
in (8) as $y_0,\ldots,y_8$: the old $\gamma_1$ is the new trunk color
$\alpha$, the old $\alpha$ is the new first-block color $\gamma_1$, and
$\beta,\gamma_2$ keep their names.  Relabel the auxiliary matching
identities by their current endpoints as well; in particular, $b_{56}$
henceforth denotes the $\beta$-identity on the current pair
$\{y_5,y_6\}$.  Relation (7), the pairwise separations in (7a), the three
differently colored identities between the final vertices, and all
conclusions above retain the same form.  Thus in either case we henceforth
have balanced data of form (2) with $\gamma_2$ missing at a trunk vertex.

We claim that some $\gamma_2$-edge joins the trunk to
$\{y_7,y_8\}$.  Suppose otherwise.  The attachment clauses for the two
distinct $\gamma_2$-colored tree identities $e_7,e_8$, together with
properness, force them to match $\{y_5,y_6\}$ to $\{y_7,y_8\}$, one at
each vertex of both pairs.  Relabel within the pairs so that
$e_7=y_5y_7$ and $e_8=y_6y_8$.  Along with the $\beta$-identity
$b_{56}$ on $y_5y_6$ and the distinct $\beta$-identity $g_\beta$ on
$y_7y_8$, these four individual identities
give the incidence cycle
\[
 y_5,b_{56},y_6,e_8,y_8,g_\beta,y_7,e_7,y_5.        \tag{8a}
\]
Indeed, expanding \noderef{TwoColorIncidenceGraph} shows that every
successive pair in (8a) is adjacent: $b_{56},g_\beta$ have color $\beta$,
$e_7,e_8$ have color $\gamma_2$, and their exact endpoint pairs are the
displayed pairs.  At each of the four vertex-nodes, (8a) already exhibits
one incident identity of each of the two distinct colors.  Properness makes
each such colored identity unique.  Equivalently, the two displayed
neighbors and \noderef{TwoColorIncidenceDegreeAtMostTwo} give exact degree
two at every vertex-node of (8a).  At an edge-node, \noderef{Incident}
says that its graph neighbors are precisely its incident endpoints; the
four exact endpoint equalities above and the distinct-end law in
\noderef{LooplessMultigraph} therefore exhibit two neighbors, and the same
degree bound shows that these are all of them.  Thus no adjacency of the
two-color incidence graph leaves the eight displayed nodes.  They form a
connected cycle containing $y_5$, so their set is exactly
\noderef{TwoColorComponent} through $y_5$, not merely a cycle contained in
that component.

Let $\varphi'$ be the switch of that component from
\noderef{TwoColorComponentSwap}.  The colors are distinct by (7), so
\noderef{TwoColorComponentSwapProper} makes $\varphi'$ proper.  The
component is disjoint from the trunk, hence it does not contain $e_1$ and
does not change any trunk incidence or missing set.  The definition of
the swap leaves the unique uncolored identity unchanged, verifying
\noderef{IsOneEdgeUncolored}.  All tree edges before $e_7$ keep their
colors and old Tashkinov witnesses; $e_7,e_8$ now have color $\beta$,
which is missing at the trunk vertex $y_1$.  Thus the same sequence is a
Tashkinov tree for $\varphi'$.  The numerical part of the $k$-triple and
the order-nine vertex set are unchanged, so these checks rebuild
\noderef{GlobalMaximumTriple} for $(T,e,\varphi')$.

In the switched coloring, the following is an incidence path with all six
edge identities displayed:
\[
 y_1,e_5,y_5,e_7,y_7,g_{\gamma_1},y_8,
 e_8,y_6,e_6,y_4,q,y_0.                            \tag{9}
\]
Here $e_5,e_6,g_{\gamma_1}$ retain color $\gamma_1$ and have exact endpoint
pairs $\{y_1,y_5\},\{y_4,y_6\},\{y_7,y_8\}$; the component switch gives
$e_7,e_8$ color $\beta$ without changing their endpoint pairs
$\{y_5,y_7\},\{y_6,y_8\}$; and $q$ retains color $\beta$ and endpoint pair
$\{y_4,y_0\}$.  Thus \noderef{TwoColorIncidenceGraph} gives every one of
the twelve successive adjacencies in (9).  The six identities are
distinct: $e_5,e_6,e_7,e_8$ are distinct tree identities by
\noderef{TreeSequence}; $q$ has endpoint pair $\{y_4,y_0\}$, different
from each of their endpoint pairs; and $g_{\gamma_1}$ has endpoint pair
$\{y_7,y_8\}$, different from the other five.  The seven vertices are
distinct by vertex injectivity in the same definition.

We prove that (9) exhausts its component.  At $y_1$, the color $\beta$
remains missing because the switched component (8a) is disjoint from the
trunk, while $e_5$ supplies an incident $\gamma_1$-identity.  Properness
and \noderef{MissingColors} therefore give degree exactly one in the
$\gamma_1,\beta$ incidence graph.  Similarly, $\gamma_1$ remains missing
at $y_0$, while $q$ supplies its incident $\beta$-identity, so $y_0$ also
has degree exactly one.  Each internal vertex-node
$y_5,y_7,y_8,y_6,y_4$ already has the two displayed incident identities,
one of each color; properness makes them unique, and
\noderef{TwoColorIncidenceDegreeAtMostTwo} makes its degree exactly two.
Each of the six edge-nodes has precisely its two displayed endpoint
incidences: this follows from its exact endpoint equality,
\noderef{Incident}, and the distinct-end clause of
\noderef{LooplessMultigraph}; the degree-at-most-two lemma again rules out
any further neighbor.  Hence no incidence adjacency leaves the thirteen
nodes in (9).  Since (9) is connected, its node set is exactly the full
\noderef{TwoColorComponent} through $y_1$ (and through $y_0$).

Apply \noderef{ScheideKempePathHeightBound} to the rebuilt globally
maximum triple, and call its returned simple $y_0$--$y_1$ incidence path
$P$.  We do not assume that the statement of that lemma makes $P$ span the
component.  Every node of $P$ nevertheless lies in the component through
$y_0$, since its initial subpath connects that node to $y_0$; the component
equality just proved therefore puts every node of $P$ among the thirteen
nodes in (9).  Reverse $P$ and start at $y_1$.  The degree-one condition at
$y_1$ forces its first edge-node to be $e_5$.  At each following internal
node of degree two, path simplicity forbids returning to the preceding
node, so the next node is forced successively to be
\[
 y_5,e_7,y_7,g_{\gamma_1},y_8,e_8,y_6,e_6,y_4,q,y_0.
\]
Thus every such simple $y_1$--$y_0$ path traverses the same ordered
support (9), all seven multigraph vertices, and exactly twelve incidence
steps.  It follows that $\mathsf{length}/2+1=12/2+1=7$.  The same lemma
bounds this quantity by $h(G)=5$, a contradiction.  The claim follows.

Let $y_j$ be the trunk vertex missing $\gamma_2$.  If $j=0$, make no
change.  If $j>0$, temporarily apply
\noderef{ScheideTrunkRotationPreserves}.  By
\noderef{TrunkRotation}, the new first vertex is the old $y_j$, all
post-trunk identities and positions are unchanged, and only the two
trunk colors are shifted on the closing cycle.  Since (7) separates
$\gamma_2$ from both trunk colors, all $\gamma_2$ identities and its
missing status are unchanged.  We may therefore apply
\noderef{ScheideHeightFiveGammaEdges} to $\gamma_2$, using the crossing
just proved, and then apply
\noderef{ScheideHeightFiveOutsideLinks} to its two returned crossing
identities.

The first lemma gives one internal $\gamma_2$-identity joining the two
middle trunk vertices and two crossing identities at the remaining two
non-missing trunk vertices.  Properness makes these exactly all
$\gamma_2$ incidences on the trunk.  The second lemma joins their outside
ends by one identity of each temporary trunk color.  At least one outside
end is $y_7$ or $y_8$ by the preceding claim.  We already have both an
$\alpha$-identity and a $\beta$-identity between $y_7,y_8$, and the
temporary trunk colors are precisely these two colors, possibly in the
opposite order.  Properness at the known final vertex forces the new
same-colored identity to be the already named identity to the other final
vertex.  Hence the two outside ends are exactly $y_7,y_8$.

Undoing the temporary rotation restores only the trunk-cycle colors and
positions; it does not alter these $\gamma_2$ identities.  Each of
$y_7,y_8$ is already incident with its tree identity $e_7$ or $e_8$ of
color $\gamma_2$.  Properness identifies those tree identities with the
two crossing identities just found.  Consequently both $e_7,e_8$ attach
to the five-vertex trunk, one at each of $y_7,y_8$.                \tag{10}

We finally build \noderef{ScheideStructuralConfiguration}.  Take $p=5$,
$s=2$, let $\Gamma=(\gamma_1,\gamma_2)$ in that order, use
$y_0,\ldots,y_4$ as the base vertices, $e_1,\ldots,e_4$ as the four base
edge identities, and use the ordered blocks
\[
 (e_5,y_5,e_6,y_6),\qquad(e_7,y_7,e_8,y_8).
\]
Equation (1) and the injective image calculation above give the required
length equation.  Equation (2) gives every exact vertex and edge
position.  Equations (6) and (10) give every block-to-base endpoint
field, and (3) gives the block colors.  The colors are injective by (7).
The first block color is missing at $y_0$, and the second at $y_j$, so
both block colors have base witnesses.  Every colored nonfirst base edge
has color $\alpha$ or $\beta$ along the normal closing path.  The two
inequalities in (7a) show that neither base color is $\gamma_1$, and (7)
shows that neither is $\gamma_2$.  Hence every such base edge avoids both
block colors.  The assumed
incidence condition on every graph vertex supplies the no-isolated field,
and the current triple is globally maximum by construction.

It remains only to verify the exit field.  By \noderef{Degree} and
\noderef{MaximumDegree}, the $k$-triple inequality gives
$d_G(v)\leq\Delta(G)\leq k-1$.  Properness makes the colors on the colored
identities incident with $v$ pairwise distinct.  By
\noderef{IsOneEdgeUncolored}, a non-end of $e$ has all $d_G(v)$ incident
identities colored and therefore misses at least $k-d_G(v)\geq1$ color;
an end of $e$ has only $d_G(v)-1$ colored incident identities and misses
at least $k-d_G(v)+1\geq2$ colors.  The exact first-edge attachment in
\noderef{TreeSequence} says that the two ends of $e$ are $y_0,y_1$.

Elementarity of $V(T)$ from \noderef{ScheideGlobalMaximumOrder}, using
\noderef{ColoredVertexSetElementary}, makes missing-color sets at distinct
tree vertices disjoint.  The four colors
used by the colored tree identities are exactly
\[
 \noderef{TreeUsedColors}(\varphi,T)
   =\{\alpha,\beta,\gamma_1,\gamma_2\};              \tag{11}
\]
their exhibited missing witnesses are respectively
$y_0,y_1,y_0,y_j$.  If $v\in V(T)\setminus\{y_0,y_j\}$ is not an end of
$e$, then it is distinct from all four witnesses and its guaranteed
missing color is outside (11).  If it is an end of $e$, it must be
$y_1$; among the four used colors only $\beta$ can be missing there, and
the count of at least two missing colors supplies another color outside
(11).  In either case $v$ misses a color which is free in the sense of
\noderef{TreeFreeColors}.

By \noderef{ScheideExitVertex}, every exit is a tree vertex.  But
\noderef{ScheideExitNoFree} says that an exit vertex has no missing free
color.  Hence every vertex in \noderef{ScheideExitVertices} lies in
$\{y_0,y_j\}$.  The first of these misses $\gamma_1$ and the second
misses $\gamma_2$, so every exit misses a member of $\Gamma$.  This is
the required exit-coverage field.  All fields of the structural
configuration have now been checked for the current edge identity,
coloring, and ordered tree-sequence, which proves the asserted existence.
\end{proof}
